\newpage
\section{仿射变换在圆锥曲线中的应用}
\label{sec:quadratic-affine-applications}

已经知道椭圆可以由圆经过伸缩得到。反过来，遇到椭圆上的问题，也可以先将它变成圆，利用圆的性质求解，再把结论还原。这样做是否方便，取决于题目中的关系怎样随变换改变。中点、分比、平行及面积比可以直接传递；长度与夹角则须重新表达。本节从这些关系出发，研究怎样选择变换并使用它。

\subsection{化椭圆为圆}
\label{subsec:quadratic-choosing-affine-transformations}

对椭圆
\begin{equation*}
    E:\quad \frac{(x-h)^2}{a^2}+\frac{(y-k)^2}{b^2}=1 ,
\end{equation*}
作变换
\begin{equation}\label{eq:quadratic-ellipse-normalization}
    u=\frac{x-h}{a} ,\qquad v=\frac{y-k}{b} ,
\end{equation}
其像为单位圆\(u^2+v^2=1\)。若原图中的点、直线分别记为\(P,l\)，就用\(P',l'\)表示它们的像。还原时使用\(x=au+h\)、\(y=bv+k\)，两个方向的面积倍率分别为\(1/(ab)\)、\(ab\)。

例如，直线\(Ax+By+C=0\)的像为
\begin{equation*}
    Aa\,u+Bb\,v+Ah+Bk+C=0 .
\end{equation*}
如果点\(M\)是弦\(PQ\)的中点，则\(M'\)是圆上弦\(P'Q'\)的中点；若\(\overrightarrow{PM}=t\overrightarrow{PQ}\)，变换后仍有\(\overrightarrow{P'M'}=t\overrightarrow{P'Q'}\)。因而图形的关联关系与共线点分比可以在圆上研究。

圆化也可以直接从配方中得到。对
\begin{equation*}
    2x^2+2xy+y^2=x^2+(x+y)^2 ,
\end{equation*}
取\(u=x,v=x+y\)即可。这个变换包含剪切，仍保持直线、平行、中点和分比，面积倍率为\(1\)。因此不必总是先求出椭圆的长短轴，才着手处理仿射关系。

\begin{samepage}
\begin{example}\label{ex:quadratic-shear-midpoint-locus}
设\(E:2x^2+2xy+y^2=1\)，\(A=(1,-1)\)。点\(B\)在\(E\)上运动且\(B\ne A\)，\(M\)是\(AB\)的中点。选择变换将\(E\)化为圆，求\(M\)的轨迹，并将所得结论还原为原坐标中的方程。
\end{example}
\end{samepage}

这里选择变换的根据是方程与题设同时变得简单。椭圆变成圆以后，原题中的定点、直线及向量关系也须一同变换，不能只换曲线的方程而沿用原来的条件。

\begin{exercise}
\mbox{}\par\nobreak

\begin{samepage}
（1）将\((x-1)^2/9+(y+2)^2/4=1\)化为单位圆，求点\((4,-2)\)及直线\(2x-3y=5\)的像。求原图中三角形面积与像三角形面积的关系。
\par
\end{samepage}

\begin{samepage}
（2）将\(x^2+4xy+5y^2=1\)化为单位圆，写出逆变换，并求该椭圆的面积。
\par
\end{samepage}
\begin{samepage}
（3）对椭圆\(2(x-1)^2+2(x-1)(y+1)+(y+1)^2=1\)，选择变换化为单位圆，求中心及面积，并写出直线\(x+y=0\)的像。
\par
\end{samepage}
\end{exercise}

\subsection{中点、平行弦与分比}
\label{subsec:quadratic-affine-chords-and-ratios}

圆的垂径关系是弦中点问题的一个直接依据。在\(E:x^2/a^2+y^2/b^2=1\)中，设一组平行弦的方向为\(d=(r,s)\)，其中点为\(M(x,y)\)。圆化后，弦方向为\(d'=(r/a,s/b)\)，中点为\(M'=(x/a,y/b)\)。圆心到弦中点的向量与弦方向垂直，故
\begin{equation}\label{eq:quadratic-affine-midpoint-diameter}
    \frac{rx}{a^2}+\frac{sy}{b^2}=0 .
\end{equation}
这正是4.3得到的直径。圆上的垂径关系还原为椭圆上的共轭关系。

\begin{figure}[htbp]
    \centering
    \includegraphics[width=0.94\linewidth]{affine-chord-correspondence.pdf}
    \caption{弦中点随伸缩对应，圆中的垂径关系还原为椭圆的共轭关系}
    \label{fig:quadratic-affine-chord-correspondence}
\end{figure}

圆中直径所对的圆周角也可以这样传递。设\(A,B\)是椭圆的一条直径的端点，\(P\)是椭圆上异于\(A,B\)的点。圆化后\(P'A'\perp P'B'\)，所以在原坐标中有
\begin{equation}\label{eq:quadratic-affine-diameter-angle}
    \frac{(x_A-x_P)(x_B-x_P)}{a^2}
    +\frac{(y_A-y_P)(y_B-y_P)}{b^2}=0 .
\end{equation}
因此\(PA,PB\)的方向相互共轭；两条直线的斜率均存在时，斜率之积为\(-b^2/a^2\)。

若\(A\)固定，\(B\)在椭圆上运动，其中点满足\(\overrightarrow{OM}=\frac12\overrightarrow{OA}+\frac12\overrightarrow{OB}\)。故\(M\)的轨迹是原椭圆以\(A\)为中心缩小一半所得的像。对过右顶点\(A=(a,0)\)的弦，其中点轨迹的方程为
\begin{equation}\label{eq:quadratic-vertex-chord-midpoints}
    \frac{(x-a/2)^2}{a^2/4}+\frac{y^2}{b^2/4}=1 .
\end{equation}
若要求弦的两端不同，则轨迹中去掉\(A\)。在圆上，这条轨迹也可由\(OM'\perp A'M'\)认识；

更一般地，若\(\overrightarrow{AM}=t\overrightarrow{AB}\)，其中\(t\ne0\)固定，则\(M=A+t(B-A)\)。固定点与动点的线性组合，直接给出了轨迹的平移与伸缩。这里先看点之间的关系，往往比先设动直线更方便。

\begin{samepage}
\begin{example}
在椭圆\(x^2/9+y^2/4=1\)中，求与\(y=x\)平行的弦的中点所在直线，并确定中点的实际取值范围。若中点给定为\(M(1,-4/9)\)，求相应弦所在直线。
\end{example}
\end{samepage}

\begin{samepage}
\begin{example}
设\(A,B\)是椭圆\(x^2/9+y^2/4=1\)上一条固定弦的两端点，\(A=(3,0)\)、\(B=(0,2)\)，\(P\)在椭圆上运动。求\(\triangle ABP\)的重心\(G\)的轨迹。若要求三角形非退化，轨迹应怎样取？
\end{example}
\end{samepage}

\begin{samepage}
\begin{example}\label{ex:quadratic-affine-hidden-midpoint}
设\(P=(1,1/2)\)，过\(P\)的两条直线分别交椭圆\(x^2/4+y^2=1\)于\(A,C\)和\(B,D\)四个不同的点，且\(AB\parallel CD\)。证明直线\(AB\)的方向固定，求其斜率；再证明过\(P\)且平行于\(AB\)的弦被\(P\)平分。
\end{example}
\end{samepage}

圆化后，平行弦与相交的连线构成圆内接梯形。利用它的对称性，可以先发现题设没有直接给出的中点关系，再接回垂径关系。

\begin{exercise}
\mbox{}\par\nobreak

\begin{samepage}
（1）求椭圆\(x^2/16+y^2/9=1\)中与向量\((2,1)\)平行的弦的中点轨迹。
\par
\end{samepage}

\begin{samepage}
（2）设\(A=(2,0)\)，\(P\)在椭圆\(x^2/4+y^2=1\)上运动，且\(\overrightarrow{AM}=2\overrightarrow{AP}\)。求点\(M\)的轨迹，再将系数\(2\)改为\(-1\)比较。
\par
\end{samepage}

\begin{samepage}
（3）椭圆\(x^2/9+y^2/4=1\)上两点\(A,B\)满足\(\overrightarrow{PA}=2\overrightarrow{BP}\)，其中\(P=(0,1)\)。求弦\(AB\)所在直线，并用圆化后的几何关系解释结果。
\par
\end{samepage}
\begin{samepage}
（4）设\(A=(-3,0)\)、\(B=(3,0)\)，\(P\)在椭圆\(x^2/9+y^2/4=1\)上且异于\(A,B\)。证明\(k_{PA}k_{PB}=-4/9\)。若将\(AB\)换成椭圆的任一直径，结论是否仍成立？用圆中的直径与圆周角解释。
\par
\end{samepage}
\begin{samepage}
（5）设\(A,B\)是椭圆\(x^2/9+y^2/4=1\)的一对直径端点，\(P\)在椭圆上且异于\(A,B\)。证明\(\triangle PAB\)的重心轨迹仍为椭圆，求其方程并说明应去掉的点。
\par
\end{samepage}
\end{exercise}

\subsection{圆幂关系的应用}
\label{subsec:quadratic-affine-power-relations}

对单位圆，过点\(P'\)作一条割线，与圆交于\(A',B'\)。在直线上选定正方向，以有向长度记\(P'A'\)、\(P'B'\)，则
\begin{equation}\label{eq:quadratic-circle-power-signed}
    \overline{P'A'}\,\overline{P'B'}=|O'P'|^2-1 .
\end{equation}
右端称为点\(P'\)对单位圆的{\heiti{}幂}。事实上，设割线的单位方向为\(d'\)，交点写成\(P'+t d'\)，代入圆方程得
\begin{equation*}
    t^2+2(P'\cdot d')t+|P'|^2-1=0 .
\end{equation*}
两个根正是这两段有向长度，由根积即得上式。对半径为\(R\)的圆，同理将右端的\(1\)改为\(R^2\)。这也给出圆内相交弦、圆外两割线的乘积关系。

现在把椭圆\(E:x^2/a^2+y^2/b^2=1\)化为单位圆。若原割线的单位方向为\(d=(r,s)\)，则任一沿该方向的位移\(t d\)的像长为
\begin{equation*}
    |t|\sqrt{\frac{r^2}{a^2}+\frac{s^2}{b^2}} .
\end{equation*}
同一直线上的所有线段具有相同的伸缩倍率，故它们的长度比保持不变。乘积则带有这个倍率的平方。记\(q(x,y)=x^2/a^2+y^2/b^2\)，由圆幂关系得到
\begin{equation}\label{eq:quadratic-ellipse-directional-power}
    \overline{PA}\,\overline{PB}
    \left(\frac{r^2}{a^2}+\frac{s^2}{b^2}\right)=q(P)-1 .
\end{equation}
因此，对于固定点\(P\)，圆化后的乘积与割线方向无关；原图中的欧氏长度乘积却一般随方向改变。比较两条不同方向的割线时，方向倍率正是需要保留的因素。

\begin{samepage}
\begin{example}\label{ex:quadratic-affine-focal-power}
（2014年安徽卷理科第14题）设\(F_1,F_2\)分别是椭圆\(E:x^2+y^2/b^2=1\)（\(0<b<1\)）的左、右焦点。过\(F_1\)的直线交椭圆于\(A,B\)，且\(|AF_1|=3|BF_1|\)、\(AF_2\perp x\)轴。求椭圆方程。
\end{example}
\end{samepage}

圆化时，原焦点成为圆内两个指定点，它们不再具有焦点的意义。但题目中的共线分比仍保持，指定的轴向关系也可直接转写。这样便能在圆中使用圆幂，而不必重新求出割线与椭圆的两个交点。

\begin{exercise}
\mbox{}\par\nobreak

\begin{samepage}
（1）过\(P=(1,0)\)的直线交椭圆\(x^2/9+y^2/4=1\)于\(A,B\)。分别在直线水平、竖直时求\(|PA|\,|PB|\)，再说明这一乘积为什么不是常数。
\par
\end{samepage}

\begin{samepage}
（2）在上题中，\(A=(3,0)\)时，求另一交点\(B\)。再设\(A=(0,2)\)，利用圆化及圆幂求\(|PA|/|PB|\)。
\par
\end{samepage}

\begin{samepage}
（3）设\(P\)在椭圆内部，证明过\(P\)的弦\(AB\)满足
\begin{equation*}
    \min(a^2,b^2)\bigl(1-q(P)\bigr)\le |PA|\,|PB|
    \le\max(a^2,b^2)\bigl(1-q(P)\bigr) ,
\end{equation*}
并指出等号对应的方向。
\par
\end{samepage}
\begin{samepage}
（4）过\(M=(1,0)\)的两条不同直线分别交椭圆\(x^2/9+y^2/4=1\)于\(A,B\)及\(C,D\)，斜率分别为\(k_1,k_2\)。若\(|MA|\,|MB|=|MC|\,|MD|\)，证明\(k_1+k_2=0\)。
\par
\end{samepage}

\begin{samepage}
（5）设\(P=(1,0)\)，过\(P\)的直线交椭圆\(x^2/9+y^2/4=1\)于\(A,B\)。过中心\(O\)作\(OC\parallel AB\)，其中\(C\)在椭圆上。证明\(|PA|\,|PB|/|OC|^2\)为定值，并求此值。
\par
\end{samepage}
\end{exercise}

\subsection{面积与面积比}
\label{subsec:quadratic-affine-area-and-extrema}

椭圆化为单位圆后，所有图形的面积都乘以同一常数\(1/(ab)\)。因此可以在圆上寻找面积最大或最小的图形，再把面积与取得条件同时还原；两个面积的比值则无需换算。

椭圆内接平行四边形的像是圆内接平行四边形，其对角既相等又互补，故为矩形，中心是圆心。还原后，原平行四边形的中心为椭圆中心。设像矩形的边长为\(r,s\)，其对角线长为\(2\)，故\(r^2+s^2=4\)，于是
\begin{equation*}
    rs\le\frac{r^2+s^2}{2}=2 .
\end{equation*}
当它是正方形时取等号，故原椭圆内接平行四边形的最大面积为\(2ab\)。最大面积的图形是圆内接正方形的仿射像，一般不是原平面中的正方形。

若以两邻接边方向看圆上的正方形，它们相互垂直；还原后，这两个方向相互共轭。以两条对角线看，它们也给出一对共轭直径。这将圆的面积最值与椭圆的共轭关系联系起来。

\begin{samepage}
\noindent\begin{minipage}{\linewidth}
\begin{example}\label{ex:quadratic-affine-area-ratio}
（2013年山东卷文科第22题）椭圆\(C\)的中心在原点，焦点在\(x\)轴上，短轴长为\(2\)，离心率为\(\sqrt2/2\)。

\begin{samepage}
（1）求椭圆方程。
\par
\end{samepage}

\begin{samepage}
（2）\(A,B\)为椭圆上满足\(S_{\triangle AOB}=\sqrt6/4\)的任意两点，\(E\)为\(AB\)的中点，射线\(OE\)交椭圆于\(P\)。设\(\overrightarrow{OP}=t\overrightarrow{OE}\)，求\(t\)。
\par
\end{samepage}
\end{example}
\end{minipage}
\end{samepage}

\begin{samepage}
对面积与中点共同出现的问题，可以在圆化后先作圆心到弦中点的连线。它既给出直角三角形，又把面积条件与中点到圆心的距离联系起来，最后利用分比保持还原。
\par
\end{samepage}

\begin{samepage}
\begin{example}
求椭圆\(x^2/9+y^2/4=1\)内接平行四边形的最大面积，分别给出一个边平行于坐标轴和一个边不平行于坐标轴的取等图形。判断所有取得最大面积的图形是否都是矩形。
\end{example}
\end{samepage}

\begin{samepage}
\begin{example}\label{ex:quadratic-affine-fixed-point-area}
过定点\(T=(0,-2)\)的直线与椭圆\(x^2/4+y^2=1\)交于不同的两点\(A,B\)。求\(S_{\triangle OAB}\)的最大值及取得最大值时直线\(AB\)的方程，其中\(O\)为原点。
\end{example}
\end{samepage}

过定点的条件在圆化后仍须保留。圆内两条半径所成三角形的面积有统一上界，但还须检查：取得这个上界的弦，能否经过指定点？最后将取等弦的方向一同还原。

\begin{exercise}
\mbox{}\par\nobreak

\begin{samepage}
（1）设\(A=(-3,0)\)、\(B=(3,0)\)，点\(P\)在椭圆\(x^2/9+y^2/4=1\)上运动。求\(S_{\triangle ABP}\)的最大值，比较直接用高与化椭为圆两种处理。
\par
\end{samepage}

\begin{samepage}
（2）求椭圆\(x^2/a^2+y^2/b^2=1\)中\(S_{\triangle OAB}\)的最大值，并用\(OA,OB\)的共轭关系说明取等位置。
\par
\end{samepage}

\begin{samepage}
（3）设\(A=(a,0)\)、\(B=(0,b)\)，点\(P\)在椭圆\(x^2/a^2+y^2/b^2=1\)上运动。求\(S_{\triangle ABP}\)的最大值及对应点\(P\)。
\par
\end{samepage}

\begin{samepage}
（4）椭圆\(x^2/t^2+y^2/(4-t^2)=1\)，其中\(0<t<2\)。求其内接平行四边形面积的最大可能值，说明为什么不能只求单位圆内接图形的面积而忽略变换倍率。
\par
\end{samepage}
\begin{samepage}
（5）点\(P\)在椭圆\(x^2/9+y^2/4=1\)上，过\(P\)且不过原点\(O\)的直线交椭圆\(x^2/36+y^2/16=1\)于\(A,B\)，射线\(PO\)交外椭圆于\(Q\)。求\(S_{\triangle ABQ}/S_{\triangle ABO}\)及\(S_{\triangle ABQ}\)的最大值。
\par
\end{samepage}

\begin{samepage}
（6）证明椭圆的任意一对共轭直径将其内部区域分成面积相等的四部分，并求\(x^2/9+y^2/4=1\)中每部分的面积。
\par
\end{samepage}

\begin{samepage}
（7）椭圆\(x^2/9+y^2/4=1\)上三点\(A,B,C\)满足\(\overrightarrow{OA}+\overrightarrow{OB}+\overrightarrow{OC}=0\)。证明\(\triangle ABC\)的面积为定值，并求此值。圆化后，这个向量条件对三角形的形状有什么要求？
\par
\end{samepage}

\begin{samepage}
（8）若不要求椭圆\(x^2/a^2+y^2/b^2=1\)的内接四边形为平行四边形，其面积能否超过\(2ab\)？证明你的判断，并确定取得最大值的条件。
\par
\end{samepage}
\end{exercise}

\newpage
\subsection{距离与夹角条件的转化}
\label{subsec:quadratic-affine-metric-conditions}

圆化不保持原长度，但有时仍能使距离问题简化。先研究椭圆\(x^2/a^2+y^2/b^2=1\)上的线性式\(Ax+By\)，其中\(A,B\)不全为零。令\(x=au,y=bv\)，则\(Ax+By=(Aa,Bb)\cdot(u,v)\)。
由于\(u^2+v^2=1\)，由数量积得
\begin{equation}\label{eq:quadratic-ellipse-linear-range}
    -\sqrt{A^2a^2+B^2b^2}\le Ax+By\le\sqrt{A^2a^2+B^2b^2} .
\end{equation}
当\((u,v)\)与\((Aa,Bb)\)同向或反向时分别取到两端。原图中的对应点为\((x,y)=\pm(Aa^2,Bb^2)/\sqrt{A^2a^2+B^2b^2}\)。
椭圆内部的点也满足上述不等式，所以区域上线性式的两端值仍在边界取得。

设直线\(l:Ax+By+C=0\)，其中\((A,B)\ne(0,0)\)。距离为\(d(P,l)=|Ax+By+C|/\sqrt{A^2+B^2}\)。
记\(R=\sqrt{A^2a^2+B^2b^2}\)，分子是区间\([C-R,C+R]\)中数的绝对值，故
\begin{equation}\label{eq:quadratic-ellipse-line-distance-range}
    d_{\max}=\frac{|C|+R}{\sqrt{A^2+B^2}} ,\qquad
    d_{\min}=\frac{\max\{|C|-R,0\}}{\sqrt{A^2+B^2}} .
\end{equation}
这里的分母来自原直角坐标中的距离公式。不能把像圆到像直线的距离不加换算地当作原来的距离。

\begin{samepage}
\begin{example}\label{ex:quadratic-affine-line-distance}
（据2019年全国I卷理科第22题改编）曲线\(C\)为椭圆\(x^2+y^2/4=1\)去掉点\((-1,0)\)所得的点集，直线\(l:2x+\sqrt3y+11=0\)。求\(C\)上的点到\(l\)距离的最小值，并确定取得最小值的点。
\end{example}
\end{samepage}

对两点间距离及夹角，需要追踪位移。若原位移的圆化表示为\(w=(\Delta u,\Delta v)\)，则原位移为\((a\Delta u,b\Delta v)\)，其平方长度为
\begin{equation}\label{eq:quadratic-restored-length}
    a^2(\Delta u)^2+b^2(\Delta v)^2 .
\end{equation}
若另一个位移表示为\(z=(\delta u,\delta v)\)，则原数量积为
\begin{equation}\label{eq:quadratic-restored-dot-product}
    a^2\Delta u\,\delta u+b^2\Delta v\,\delta v .
\end{equation}
所以等长可比较式\eqref{eq:quadratic-restored-length}，垂直则令式\eqref{eq:quadratic-restored-dot-product}为零；夹角由这个数量积除以两段原长度的乘积得到。圆的简单方程与这些度量表达可以同时使用。

\begin{samepage}
\begin{example}
在椭圆\(x^2/9+y^2/4=1\)上，\(A=(3,0)\)，点\(B\ne A\)满足\(OA\perp AB\)是否可能？若将条件改为\(OB\perp AB\)，求所有这样的点\(B\)。分别用原坐标与圆化坐标处理，比较需要转化的条件。
\end{example}
\end{samepage}

\begin{exercise}
\mbox{}\par\nobreak

\begin{samepage}
（1）求\(x^2/9+y^2/4\le1\)中\(2x-y\)的最大值、最小值及取等点。
\par
\end{samepage}

\begin{samepage}
（2）求椭圆\(x^2/9+y^2/4=1\)到直线\(x+y=5\)的最近距离与最远距离。若直线改为\(x+y=1\)，最近距离怎样变化？
\par
\end{samepage}

\begin{samepage}
（3）（2013年上海卷秋季文科第18题）设\(\Omega_n\)为\(x^2/4+ny^2/(4n+1)\le1\)确定的区域，\(M_n\)为\(x+y\)在\(\Omega_n\)上的最大值。求\(M_n\)及\(\lim_{n\to\infty}M_n\)。
\par
\end{samepage}
\begin{samepage}
（4）椭圆\(x^2/9+y^2/4=1\)的两条共轭半径\(OA,OB\)等长。求它们的长度及夹角余弦的绝对值，说明圆化后应同时使用哪些关系。
\par
\end{samepage}

\begin{samepage}
（5）点\(P\)在椭圆\(x^2/9+y^2/4=1\)上，\(T=(0,1)\)。求\(|PT|\)的最大值、最小值及取等点。可先圆化，将距离的平方化为一个坐标的二次式；比较本题与点到定直线距离问题的区别。
\par
\end{samepage}

\begin{samepage}
（6）椭圆\(x^2/9+y^2/4=1\)上两条共轭半径为\(OA,OB\)。证明\(|OA|^2+|OB|^2=13\)，求\(|OA|\,|OB|\)的取值范围，并指出两端值对应的方向。说明共轭关系与欧氏长度怎样在圆化后配合使用。
\par
\end{samepage}

\begin{samepage}
（7）设\(A=(-3,0)\)、\(B=(3,0)\)，点\(P\)在椭圆\(x^2/9+y^2/4=1\)上。求\(|PA|\,|PB|\)的最大值及取等点。若先将椭圆化为单位圆，能否直接将像图中的距离乘积乘以一个固定常数还原？说明理由。
\par
\end{samepage}

\begin{samepage}
（8）设\(P=(x,y)\)在椭圆\(x^2/9+y^2/4=1\)上。求\(|x+y|+|x-y|\)的最大值、最小值及取等点。先研究绝对值式与\(\max\{|x|,|y|\}\)的关系，再比较直接用原图与圆化处理的方便程度。
\par
\end{samepage}

\begin{samepage}
（9）椭圆\(x^2/9+y^2/4=1\)的内接矩形各边平行于坐标轴。求其周长的最大值及取等顶点，再求面积的最大值。取得最大周长与最大面积的矩形是否相同？联系线性式的最值解释。
\par
\end{samepage}

\begin{samepage}
（10）椭圆\(x^2/a^2+y^2/b^2=1\)（\(a>b>0\)）上，\(x+y\)的最大值为\(\sqrt{13}\)，\(2x+y\)的最大值为\(2\sqrt{10}\)。求椭圆方程，并求上述两个线性式各自取得最大值的点。
\par
\end{samepage}

\begin{samepage}
（11）设\(A,B\)是椭圆\(x^2/9+y^2/4=1\)上两点，满足\(|OA|=|OB|\)、\(\angle AOB=60^\circ\)。求\(\triangle OAB\)面积的所有可能值。利用原图中的对称性研究等长条件，再与圆化后的度量表达比较。
\par
\end{samepage}
\end{exercise}

\newpage
\subsection{综合应用与方法选择}
\label{subsec:quadratic-affine-problem-solving}

变换的作用应当体现在具体关系的简化上。若主要条件是共线、平行、中点或面积比，可以直接利用它们的保持性；若还含长度、夹角，则须看新度量表达与圆的方程能否配合。原图中已经容易处理的关系，可以留在原图中完成。

\begin{samepage}
\begin{example}\label{ex:quadratic-affine-method-choice}
（2018年天津卷文科第19题）椭圆\(x^2/a^2+y^2/b^2=1\)（\(a>b>0\)）的右顶点为\(A\)，上顶点为\(B\)，离心率为\(\sqrt5/3\)，且\(|AB|=\sqrt{13}\)。

\begin{samepage}
（1）求椭圆方程。
\par
\end{samepage}

\begin{samepage}
（2）直线\(l:y=kx\)（\(k<0\)）交椭圆于\(P,Q\)，交\(AB\)于\(M\)，点\(P,M\)均在第四象限。若\(S_{\triangle BPM}=2S_{\triangle BPQ}\)，求\(k\)。
\par
\end{samepage}
\end{example}
\end{samepage}

本题的共底边面积比与中心对称直接给出共线分比。可以圆化后研究，也可以在原图中从分比写出坐标关系，再代入椭圆方程。比较两种处理时，应看圆化是否进一步简化了这些关系。

\begin{exercise}
\mbox{}\par\nobreak

\begin{samepage}
（1）（据2012年湖北卷理科第21题）单位圆上一点\(A\)到\(x\)轴的垂足为\(D\)，点\(M\)在直线\(AD\)上，满足\(|DM|=m|DA|\)，其中\(m>0\)、\(m\ne1\)。求\(M\)的轨迹\(C\)及焦点。过原点、斜率为\(k>0\)的直线交\(C\)于\(P,Q\)，\(P\)在第一象限，\(N\)为\(P\)在\(y\)轴上的射影，\(QN\)再交\(C\)于\(H\)。是否存在\(m\)，使得任意\(k>0\)都有\(PQ\perp PH\)？比较第二问留在原坐标与圆化后还原垂直条件的处理。
\par
\end{samepage}

\begin{samepage}
（2）设\(P\)为椭圆\(x^2/9+y^2/4=1\)外一点。两条过\(P\)的直线分别交椭圆于\(A,B\)和\(C,D\)，其单位方向分别为\(d_1,d_2\)。写出\(|PA|\,|PB|\)与\(|PC|\,|PD|\)之比。两条直线方向互相垂直时，这两个乘积之和是否必为定值？
\par
\end{samepage}

\begin{samepage}
（3）设\(A=(3,0)\)，\(B=(0,2)\)，\(P\)在椭圆\(x^2/9+y^2/4=1\)上。分别求\(S_{\triangle ABP}\)、\(|OP|\)的最大值。判断这两个目标是否适合用同一种方法处理。
\par
\end{samepage}

\begin{samepage}
（4）设\(A,B\)是椭圆\(x^2/a^2+y^2/b^2=1\)上的点，且\(OA\perp OB\)。求\(|OA|^{-2}+|OB|^{-2}\)，并求\(S_{\triangle OAB}\)的取值范围。比较按原图中的方向设坐标与直接化椭为圆两种处理。
\par
\end{samepage}

\begin{samepage}
（5）求椭圆\(x^2/9+y^2/4=1\)内接三角形面积的最大值，并给出一组取等顶点。取得最大面积时，原图中的三角形是否必须为等边三角形？
\par
\end{samepage}
\begin{samepage}
（6）求椭圆\(2x^2+2xy+y^2=1\)内接平行四边形面积的最大值，给出一组取等顶点，并判断该图形是否为矩形。
\par
\end{samepage}
\end{exercise}

\newpage
\subsection{仿射变换与度量*}
\label{subsec:quadratic-affine-metric-forms}

4.5.6已区分描述曲线与描述原长度的二次型。现在将本节的长度与数量积公式推广到一般仿射变换。设\(X=SU+t\)，其中\(S\)可逆。平移在位移中消去，故对于\(\Delta X=S\Delta U\)、\(\Delta Y=S\Delta V\)，有
\begin{equation}\label{eq:quadratic-pulled-back-metric}
    |\Delta X|^2=\Delta U^TG\Delta U ,\qquad
    \Delta X\cdot\Delta Y=\Delta U^TG\Delta V ,\qquad G=S^TS .
\end{equation}
对非零\(W\)，\(W^TGW=|SW|^2>0\)，因此\(G\)正定。它给出新坐标中原欧氏度量的表达。沿新坐标方向\(W\)的长度倍率为\(\sqrt{W^TGW}/|W|\)，而原图中的正交关系成为\(U^TGV=0\)。

对以中心为原点的椭圆\(X^THX=1\)，取\(t=0\)，选择\(S\)使\(S^THS=I\)，则其新方程是单位圆，原长度与夹角由\(G\)计算。4.5.6中伸缩变换的度量表达，是这一公式的特殊情形。

在二维情形，\(\det G=(\det S)^2\)。所以\(\sqrt{\det G}=|\det S|\)也是面积倍率。前面的长度、方向乘积与面积计算由此接到同一个线性变换上。

\begin{samepage}
\begin{example}
对\(2x^2+2xy+y^2=1\)，取\(u=x,v=x+y\)。写出逆变换矩阵\(S\)及度量矩阵\(G\)，分别求像圆上两个坐标方向还原后的长度倍率、夹角与面积倍率。用这一表达重新检查本节第一个中点轨迹例题。
\end{example}
\end{samepage}

\begin{exercise}
\mbox{}\par\nobreak

\begin{samepage}
（1）（2012年福建卷理科第21题）曲线\(2x^2+2xy+y^2=1\)在矩阵\(A=\left[\begin{smallmatrix}a&0\\b&1\end{smallmatrix}\right]\)（\(a>0\)）对应的变换下成为单位圆。求\(a,b\)及\((A^2)^{-1}\)，再写出圆化坐标中的原度量矩阵。
\par
\end{samepage}

\begin{samepage}
（2）证明若两个可逆矩阵\(S_1,S_2\)满足\(S_1S_1^T=S_2S_2^T\)，则\(S_2^{-1}S_1\)保持长度。解释同一椭圆的不同圆化表示之间的关系。
\par
\end{samepage}

\begin{samepage}
（3）设\(H\)正定，\(q(X)=X^THX\)。沿非零方向\(V\)的直线\(X=P+tV\)交\(q(X)=1\)于两个点，参数为\(t_1,t_2\)。证明\(t_1t_2q(V)=q(P)-1\)，并与式\eqref{eq:quadratic-ellipse-directional-power}比较。
\par
\end{samepage}
\begin{samepage}
（4）设\(S^THS=I\)、\(G=S^TS\)。证明\(H^{-1}=SS^T\)，且\(G\)与\(H^{-1}\)有相同的特征值。由此解释圆化坐标中的最大、最小长度倍率与原椭圆半轴长的关系。
\par
\end{samepage}

\begin{samepage}
（5）原平面中的单位圆在坐标变换\(X=SU\)后成为怎样的曲线？用\(G=S^TS\)写出方程，求新坐标平面中该曲线所围区域的面积，并与面积倍率核对。
\par
\end{samepage}

\begin{samepage}
（6）证明：若椭圆的每一对共轭方向都互相垂直，则该椭圆为圆。分别用曲线矩阵\(H\)与圆化后的度量矩阵\(G\)解释这一结论。
\par
\end{samepage}
\end{exercise}
